
\chapter{JSON}
\label{ch:JSON}

JSON (pronounced like the name ``Jason'') is a simple, human-readable text file
format that is very commonly used for information storage and interchange. (It
stands for ``JavaScript Object Notation,'' but don't be fooled: it doesn't have
much to do with the JavaScript language except historically. \textit{Any}
language, including Python, can read/write JSON data.)

In this chapter, we'll learn about how to read and navigate through JSON data.
First, though, a word about a couple of Python essentials I omitted from DATA
101 for brevity and simplicity.

\section{Lists and dictionaries}

\index{series@\texttt{Series} (Pandas)}
\index{ndarray@\texttt{ndarray} (NumPy)}
\index{associative array}
\index{key-value pair}
\index{index}

Last semester, we spent a lot of time learning how to use NumPy
\texttt{ndarray}s and Pandas \texttt{Series}es. These are the preferred
implementations of arrays, and associative arrays, respectively, in the Python
Data Science ecosystem. As you know, a NumPy array holds a numbered sequence of
items (indexed starting from 0), and a \texttt{Series} holds a set of key-value
pairs.

For any serious data work, you want to use those tools, because they are
lightning-fast, super-efficient, feature-rich, and optimized to the task.

\index{plain-ol' Python}

As it turns out, plain-ol' Python has stripped-down, feature-poor, but
easy-to-type versions of these data structures as well, which are baked in to
JSON and also important to know about in their own right. They're called
\textbf{lists} and \textbf{dictionaries}.

\subsection{Lists}

If, instead of writing this:

\begin{Verbatim}[fontsize=\small,samepage=true,frame=single,framesep=3mm]
crew = np.array(["Ed", "Kelly", "Alara", "Bortus", "John", "Claire"])
\end{Verbatim}

you simply wrote this:

\begin{Verbatim}[fontsize=\small,samepage=true,frame=single,framesep=3mm]
crew = ["Ed", "Kelly", "Alara", "Bortus", "John", "Claire"]
\end{Verbatim}

\index{list}

then you get a plain-ol \textbf{list} instead of a NumPy array. It isn't nearly
as fast, nor can it store nearly as many items, and you can't do broadcasting
operations and such with it, but it does the basic job. For example:

\begin{Verbatim}[fontsize=\small,samepage=true,frame=single,framesep=3mm]
print("Captain {} initiating roll call for {} officers:".format(
    crew[0], len(crew)))
for member in crew[1:len(crew)]:
    print("{} reporting for duty, sir!".format(member))
print(crew[0])
\end{Verbatim}

\vspace{-.2in}

\begin{Verbatim}[fontsize=\small,samepage=true,frame=leftline,framesep=5mm,framerule=1mm]
Captain Ed initiating roll call for 6 officers:
Kelly reporting for duty, sir!
Alara reporting for duty, sir!
Bortus reporting for duty, sir!
John reporting for duty, sir!
Claire reporting for duty, sir!
\end{Verbatim}

Working with small lists of items is very, very similar to working with NumPy
arrays, so I'm sure you'll pick this right up.

Incidentally, we've actually used lists several times already, without knowing
it. For instance, this very line:

\begin{Verbatim}[fontsize=\small,samepage=true,frame=single,framesep=3mm]
good = np.array(["Janet", "Chidi", "Eleanor", "Michael", "Tahani"])
\end{Verbatim}

is actually first creating a \textit{list}, and then passing that list as an
argument to \texttt{np.array()}. And this code:

\begin{Verbatim}[fontsize=\small,samepage=true,frame=single,framesep=3mm]
alter_egos = pd.Series(['Hulk','Spidey','Iron Man','Thor'],
    index=['Bruce','Peter','Tony','Thor'])
\end{Verbatim}

creates two separate lists, one for the data and one for the index, and passes
those to \texttt{pd.Series()} to build a \texttt{Series} out of them. Lists are
super-lightweight and super-common.


\pagebreak

\section{JSON}

\index{JSON}

The really nice things are (1) it's simple, and easy on the eyes (unlike
binary formats, or XML), and (2) it structures information hierarchically in
list-ish, dictionary-ish, string-ish type elements that will be instantly
familiar to you. A JSON ``document'' is, in fact, \textit{essentially a Python
list or dictionary with embedded contents.}

Here's an example file in JSON format:

\small
\begin{Verbatim}[fontsize=\small,samepage=true,frame=none]
{"name": "Luke Skywalker", "rank": "Jedi", "lightsaber color": "blue",
"episodes": ["IV", "V", "VI", "VII", "VIII"], "relations": {"Leia": "sister",
"Han": "friend", "Artoo": "servant", "Rey": "student"}}
\end{Verbatim}
\normalsize

Its structure is instantly apparent: a dictionary whose values are strings,
lists, and other dictionaries, as appropriate.

Probably the best thing about JSON is that once we read and parse a JSON file
or other information source, we have a bona fide Python object that we can
treat in the ordinary way.

In Python, we have the \texttt{json} package, which you can import like any
other:

\begin{Verbatim}[fontsize=\small,samepage=true,frame=single]
    import json
\end{Verbatim}

There are four functions that matter:
\small
\begin{compactitem}
\item \texttt{json.dump}(\textit{someObject}, \textit{aFile}) - Dump the Python
object to a file in JSON text format.
\item \texttt{json.dumps}(\textit{someObject}) - Dump the Python
object to a string in JSON text format (and return it).
\item \texttt{json.load}(\textit{aFile}) - Build and return a Python
object from the contents of the JSON text file.
\item \texttt{json.loads}(\textit{someString}) - Build and return a Python object
from the JSON text string.
\end{compactitem}

\normalsize
So, for instance, to load a file \texttt{info.json} into a Python object
called ``\texttt{my\_stuff}'':

\small
\begin{Verbatim}[fontsize=\small,samepage=true,frame=single]
    with open('info.json','r') as f:
        my_stuff = json.load(f)
\end{Verbatim}

\normalsize
We could parse and access the Luke Skywalker file just like this:

\small
\begin{Verbatim}[fontsize=\small,samepage=true,frame=single]
    with open('luke.json','r') as f:
        luke = json.load(f)
    print(luke['lightsaber color'])
    blue

    print(luke['episodes'][2])
    VI
\end{Verbatim}

\normalsize
Simple as that!

\section{``Hierarchical'' vs. ``flat'' data}

You can see that the JSON file format is essentially a textual representation
of an arbitrarily nested Python object. This object can be a list or a
dictionary, and have embedded lists and dictionaries to any arbitrary level of
nesting. Here's another example:

\begin{Verbatim}[fontsize=\small,samepage=true,frame=none]
[{"prefix": "DATA", "num": 219, "creds": 3, "instructor": "Davies",
"meetings": [{"time": {"day": "TR", "time": "2pm"}, "location": {"bldg":
"Trinkle", "room": "B6"}}], "roster": [{"first": "Lady", "last": "Gaga",
"email": "lgaga@umw.edu"}, {"first": "The", "last": "Rock", "email":
"trock@umw.edu"}, ...
\end{Verbatim}

When you disentangle this, you'll realize that what you have is:

\vspace{.2in}

\large
\noindent A list, each element of which is:\\
$\bullet$ a dictionary, with key/value pairs:
        \begin{compactenum}
        \item \texttt{prefix}: a string
        \item \texttt{num}: an int
        \item \texttt{creds}: an int
        \item \texttt{instructor}: a string
        \item \texttt{meetings}: a list, each element of which is:
            \begin{compactitem}
            \item a dictionary, with key/value pairs:
                \begin{compactenum}
                \item \texttt{time}: a dictionary, with key/value pairs:
                    \begin{compactenum}
                    \item \texttt{day}: a string
                    \item \texttt{time}: a string
                    \end{compactenum}
                \item \texttt{location}: a dictionary, with key/value pairs:
                    \begin{compactenum}
                    \item \texttt{bldg}: a string
                    \item \texttt{room}: a string
                    \end{compactenum}
                \end{compactenum}
            \end{compactitem}
        \item \texttt{roster}: a list, each element of which is:
            \begin{compactitem}
            \item a dictionary, with key/value pairs:
                \begin{compactenum}
                \item \texttt{first}: a string
                \item \texttt{last}: a string
                \item \texttt{email}: a string
                \end{compactenum}
            \end{compactitem}
        \item ...and possibly other stuff...
        \end{compactenum}

\normalsize

\pagebreak

This kind of structure is called \textbf{hierarchical} data. This is in
contrast to so-called \textbf{flat} data, like a CSV:

\begin{Verbatim}[fontsize=\small,samepage=true,frame=none]
Prefix,Num,Instructor,Creds,Bldg,Room,Day,Time
DATA,219,Davies,3,Trinkle,B6,TR,2pm
ENGL,286,Rigsby,3,Combs,112,MWF,11am
CPSC,305,Finlayson,4,Trinkle,B52,MWF,8am
\end{Verbatim}

It should be clear that hierarchical data is intrinsically \textit{more
expressive} than flat data. Rather than each data record being a strict,
identical-length sequence of values, the whole nested structure communicates
order and relationships between subcomponents, like an outline.

It's not even always possible, let alone straightforward, to convert from a
hierarchical format to a flat format. Consider the JSON file above: it permits
any number of \texttt{meetings} for a single course, each of which has
embedded structure. In this case, the structure for each meeting is regular,
since it's apparently composed of a fixed number of fields with unique names.
(This won't always be the case, and if it's not, we would \textit{lose
information} in collapsing it to a flat format.) Even so, though, it's not
clear what to do with the multiple meeting times/locations for a course. Do we
make each \texttt{Bldg/Room/Day/Time} group into its own set of headings in
the flat format? This brings with it its own problems, such as:

\begin{compactenum}
\item How do we associate each group of buildings/rooms/days/times together,
since they're now all individual top-level fields? (Naming conventions seem
like the only solution, and it's awkward.)
\item How do we know the maximum number of such meetings, so as to specify the
width of the table?
\item What do we do for rows which don't have the maximum number of meetings?
Leave the others blank?
\item How do we perform standard queries like ``find all courses that have at
least one meeting in Trinkle?'' (We have to split this up into an awkward query
that separately queries each of the columns and unions the answer.)
\end{compactenum}

There are better ways to organize such data, such as the \textbf{relational
model} that has been ubiquitous in database schema design since Ted Codd
invented it in the 1960's. It allows us to ``normalize'' such data sets and
represent them in multiple, interrelated tables. We'll take a look at this
when we look at SQLite later in the course. Some irregular structures defy
even the relational model, however, which is why some newer database paradigms
(like NoSQL databases) have been developed to get around its limitations.

One big reality to grasp, however, is that at this point in time,
\textit{almost all machine learning algorithms are based on a \textbf{flat}
data structure}. There are some quite advanced, cutting-edge exceptions to
this, but really the field of mining hierarchical (and other non-flat) data
formats is still in its infancy. Lots of opportunity to make a name for
yourself here!
